\section{Contributions and Implications}
\label{sec:contributions}

The thesis makes three related contributions. Empirically, it tests the incremental predictive content of conventional ownership measures and ownership-network topology for a forward stock-level downside outcome. This extends the ownership-fragility literature, which links ownership structure and crowding to volatility and drawdowns, by evaluating the information through a fixed out-of-sample forecasting design \parencite{greenwood2011stock,brown2022crowded}. The results distinguish limited information in conventional ownership variables from the absence of a reliable incremental gain from the network representations tested here. This conclusion is specific to the variables, graph construction, models, and sample used in the thesis; it does not imply that institutional ownership or financial networks are generally uninformative.

Methodologically, the tabular comparisons change one dimension at a time. Ridge and XGBoost are compared on matched information sets, and graph models are evaluated separately as a method for learning directly from ownership relations. Point-in-time feature construction, purged chronological boundaries, fitting-sample preprocessing, and validation-stage model selection preserve the sequence in which information would have become available. The selected specification is fixed before the test sample is examined, and the held-out comparison is used to assess generalization rather than to reopen selection. This design makes it possible to distinguish small predictive gains from the effects of additional inputs or model complexity.

Economically, the thesis converts predicted drawdown into a Fragility Score computed within each quarter and evaluates it sequentially as a risk ranking, a long-only exclusion screen, and an exploratory long--short signal. The evidence is strongest at the first stage: the score consistently orders subsequent downside outcomes. Its long-only benefit is visible in the equal-weighted universe but nearly disappears under value weighting, while the long--short results depend on a short sample and simplified implementation assumptions. The appropriate implication is therefore narrower than a trading claim. The score may support risk monitoring, security review, or exclusion-based screening within a defined universe; the thesis does not establish that it should determine portfolio weights or serve as a stand-alone return signal.
